\begin{lemma}[$D$ is linear on $W^{1,2}(\gamma_n)$ and annihilates constants]
Let $W^{1,2}(\gamma_n)$ be as in \noderef{SobolevW12} and $D$ as in
\noderef{GaussD}.  Then:
\begin{enumerate}
\item for all $u, v \in W^{1,2}(\gamma_n)$, also $u + v \in W^{1,2}(\gamma_n)$ and
$D(u+v) = Du + Dv$;
\item for every $c \in \mathbb{R}$ and $u \in W^{1,2}(\gamma_n)$, also
$c\,u \in W^{1,2}(\gamma_n)$ and $D(c\,u) = c\,Du$;
\item for every $c \in \mathbb{R}$ the constant function $c\mathbf{1} = c\,h_0$ of
\noderef{HermiteL2} and \noderef{Hermite1Zero} lies in $W^{1,2}(\gamma_n)$ and
$D(c\mathbf{1}) = 0$.
\end{enumerate}
\end{lemma}
\begin{proof}
Throughout, recall from \noderef{SobolevW12} that $w \in W^{1,2}(\gamma_n)$ means
precisely that there is $\Gamma \in L^2(\gamma_n;\mathbb{R}^n)$ such that
$\sqrt2\,\Gamma$ is a weak Euclidean gradient of $w$ in the sense of
\noderef{HasGaussWeakGrad}, and from \noderef{GaussD} that $Dw$ is then that
$\Gamma$; it is unique as an element of $L^2(\gamma_n;\mathbb{R}^n)$ by
\noderef{GaussWeakGradUnique}.  Note also that the defining identity of
\noderef{HasGaussWeakGrad}, namely
\[
  \int_{\mathbb{R}^n} w\,\partial_j\varphi\, dx
    = -\int_{\mathbb{R}^n} G_j\,\varphi\, dx
  \qquad (j = 1,\ldots,n,\ \varphi \in C_c^\infty(\mathbb{R}^n)),
\]
is, for each fixed $j$ and $\varphi$, linear in the pair $(w, G)$: both sides are
linear in $(w,G)$ separately.

(1) Let $u,v \in W^{1,2}(\gamma_n)$ and put $G = \sqrt2\,Du$,
$H = \sqrt2\,Dv$, the weak Euclidean gradients of $u$ and $v$.  Then
$G + H \in L^2(\gamma_n;\mathbb{R}^n)$, and adding the two defining identities
for the same $j$ and $\varphi$ gives
\[
  \int (u+v)\,\partial_j\varphi\, dx
    = -\int (G_j + H_j)\,\varphi\, dx ,
\]
so $G+H$ is a weak Euclidean gradient of $u+v$.  Writing
$G + H = \sqrt2\,(Du + Dv)$ with $Du + Dv \in L^2(\gamma_n;\mathbb{R}^n)$, the
criterion of \noderef{SobolevW12} is met with the witness $Du+Dv$, so
$u + v \in W^{1,2}(\gamma_n)$; and since by \noderef{GaussD} the element
$D(u+v)$ is a witness too, \noderef{GaussWeakGradUnique} applied to the
function $u+v$ and the two fields $\sqrt2\,D(u+v)$ and $\sqrt2(Du+Dv)$ gives
$D(u+v) = Du + Dv$.

(2) Identical, multiplying the single defining identity for $u$ by the scalar
$c$: $c\,G = \sqrt2\,(c\,Du)$ is a weak Euclidean gradient of $c\,u$, so
$c\,u \in W^{1,2}(\gamma_n)$ with witness $c\,Du$, and
\noderef{GaussWeakGradUnique} gives $D(c\,u) = c\,Du$.

(3) By \noderef{Hermite1Zero} and \noderef{HermiteL2}, $h_0$ is the constant
function $1$, so $c\,h_0$ is the constant function $c$; it lies in
$L^2(\gamma_n)$ because $\gamma_n$ is a probability measure
(\noderef{GaussIsProb}).  We claim the zero field is a weak Euclidean gradient
of it.  Indeed, let $j \in \{1,\ldots,n\}$ and $\varphi \in C_c^\infty(\mathbb{R}^n)$.
Then $\partial_j \varphi$ is smooth with compact support, and integrating in the
$x_j$ variable first (Fubini applies, the integrand being continuous with compact
support, hence absolutely integrable) the fundamental theorem of calculus gives
\[
  \int_{\mathbb{R}} \partial_j \varphi(x)\, dx_j = 0
  \qquad\text{for every fixed value of the other coordinates,}
\]
because $\varphi$ vanishes outside a compact set, so the inner integral is
$\varphi\big|_{x_j = +\infty} - \varphi\big|_{x_j=-\infty} = 0$.  Hence
$\int_{\mathbb{R}^n}\partial_j\varphi\, dx = 0$ and therefore
\[
  \int_{\mathbb{R}^n} c\,\partial_j\varphi\, dx = 0 = -\int_{\mathbb{R}^n} 0\cdot\varphi\, dx ,
\]
which is the defining identity of \noderef{HasGaussWeakGrad} for the function
$c\mathbf{1}$ and the field $G = 0$.  Since $0 = \sqrt2\cdot 0$ with
$0 \in L^2(\gamma_n;\mathbb{R}^n)$, the criterion of \noderef{SobolevW12} holds
with witness $0$, so $c\mathbf{1} \in W^{1,2}(\gamma_n)$; and
\noderef{GaussWeakGradUnique}, applied to $c\mathbf{1}$ and the two fields
$\sqrt2\,D(c\mathbf{1})$ and $0$, gives $D(c\mathbf{1}) = 0$.
\end{proof}
